%% LyX 2.5.1 created this file.  For more info, see https://www.lyx.org/.
%% Do not edit unless you really know what you are doing.
\documentclass[10pt,hebrew,english]{article}
\usepackage{amstext}
\usepackage{amssymb}
\usepackage{fontspec}
\setmainfont[Mapping=tex-text]{Times New Roman}
\usepackage{array}
\usepackage{longtable}
\usepackage{float}
\usepackage{calc}
\usepackage{varwidth}
\usepackage{geometry}
\geometry{verbose,tmargin=1cm,bmargin=1cm,lmargin=1cm,rmargin=1cm}
\usepackage{esint}

\makeatletter

%%%%%%%%%%%%%%%%%%%%%%%%%%%%%% LyX specific LaTeX commands.
\newcommand{\lyxmathsym}[1]{\ifmmode\begingroup\def\b@ld{bold}
  \text{\ifx\math@version\b@ld\bfseries\fi#1}\endgroup\else#1\fi}

%% Because html converters don't know tabularnewline
\providecommand{\tabularnewline}{\\}
%% Variable width box for table cells
\newenvironment{cellvarwidth}[1][t]
    {\begin{varwidth}[#1]{\linewidth}}
    {\@finalstrut\@arstrutbox\end{varwidth}}

%%%%%%%%%%%%%%%%%%%%%%%%%%%%%% User specified LaTeX commands.
% הגדרת חבילת array (נחוצה לרווח אנכי)
% אילוץ שבירת שורות אוטומטית ויישור לימין בכל מיניפייג' במסמך
\usepackage{etoolbox}
\usepackage{ragged2e}
\AtBeginEnvironment{minipage}{\RaggedLeft}

\changefontsizes{8pt} % שנה את המספר 12 לגודל הרצוי לך

\usepackage{array}

% הגדרת רווח אנכי (גובה השורות) - מגדיל את המרווח מעל ומתחת לטקסט
\renewcommand{\arraystretch}{2}

% הגדרת רווח אופקי (בין העמודות) - ברירת המחדל היא 6pt
\setlength{\tabcolsep}{7pt}

% אופציונלי: הוספת רווח נוסף ספציפית מעל תווים גבוהים (כמו שברים או סוגריים)
\setlength{\extrarowheight}{0pt}

% חבילה לשליטה בכותרות
\usepackage{caption}

\makeatother

\usepackage{polyglossia}
\setdefaultlanguage[variant=american]{english}
\setotherlanguage{hebrew}
\begin{document}
\begin{hebrew}%
{\fboxsep 1pt\fbox{\begin{minipage}[t]{0.5\columnwidth}%

\section*{קבועים}

\begin{longtable}[c]{r|r}
קבוע פלאנק & $h=6.626\times10^{-34}\left[\mathrm{J\cdot s}\right]$\tabularnewline
\hline 
קבוע פלאנק המצומצם & \textenglish[variant=american]{$\hslash=\frac{h}{2\pi}=1.055\times10^{-34}\left[\mathrm{J\cdot s}\right]$}\tabularnewline
\hline 
קבוע בולצמן: & \textenglish[variant=american]{$k_{B}=1.38\times10^{-23}\left[\frac{\text{J}}{\text{K}}\right]$}\tabularnewline
\hline 
מהירות האור בריק & $c=3\times10^{8}\left[\frac{\text{m}}{\text{s}}\right]$\tabularnewline
\hline 
קבוע \textbf{סטפן-בולצמן} & $\sigma=5.672\times10^{-8}\left[\frac{\text{W}}{\text{m}^{2}\text{K}^{4}}\right]$\tabularnewline
\hline 
קבוע וין & $b=2.9\times10^{-3}\left[\text{ m}\cdot\text{K}\right]$\tabularnewline
\hline 
מסת האלקטרון & \textenglish[variant=american]{$m_{e}=9.109\times10^{-31}\text{\ensuremath{\left[\mathrm{kg}\right]}}$}\tabularnewline
\hline 
מטען האלקטרון & \textenglish[variant=american]{$e=1.602\times10^{-19}\text{\ensuremath{\left[\mathrm{C}\right]}}$}\tabularnewline
\hline 
רדיוס בוהר & \textenglish[variant=american]{$a_{0}=0.529\text{\AA}$}\tabularnewline
\hline 
אנרגיית רמת היסוד של המימן & \textenglish[variant=american]{$E_{0}=13.6\text{\ensuremath{\left[\mathrm{eV}\right]}}$}\tabularnewline
\hline 
מספר אבוגדרו & \textenglish[variant=american]{$N_{A}=6.022\times10^{23}\left[\text{mol}^{-1}\right]$}\tabularnewline
\end{longtable}%
\end{minipage}}}~%
{\fboxsep 1pt\fbox{\begin{minipage}[t]{0.5\columnwidth}%

\section*{קבועים עבור יחידות $\mathrm{eV}$}

\begin{longtable}[c]{r|r}
קבוע פלאנק & \textenglish[variant=american]{$h=4.136\times10^{-15}\left[\text{ eV}\cdot\text{s}\right]$}\tabularnewline
\hline 
קבוע פלאנק המצומצם & \textenglish[variant=american]{$\hslash=6.582\times10^{-16}\left[\text{ eV}\cdot\text{s}\right]$}\tabularnewline
\hline 
קבוע בולצמן: & \textenglish[variant=american]{$k_{B}=8.617\times10^{-5}\left[\frac{\text{eV}}{\text{K}}\right]$}\tabularnewline
\hline 
\textenglish[variant=american]{} & \textenglish[variant=american]{}\tabularnewline
\hline 
\textenglish[variant=american]{} & \textenglish[variant=american]{}\tabularnewline
\hline 
\textenglish[variant=american]{} & \textenglish[variant=american]{}\tabularnewline
\hline 
\textenglish[variant=american]{} & \textenglish[variant=american]{}\tabularnewline
\end{longtable}%
\end{minipage}}}

\fbox{\begin{minipage}[t]{0.5\columnwidth}%

\section*{חוקים בסיסיים}

\begin{longtable}[c]{r|r}
חוק וין & \textenglish[variant=american]{$\lambda_{max}=\frac{b}{T}$$\left[\mathrm{m}\right]$}\tabularnewline
\hline 
חוק סטפן-בולצמן & \textenglish[variant=american]{$I=\sigma T^{4}\left[\mathrm{\frac{W}{m^{2}}}\right]$}\tabularnewline
\hline 
הקרנה & \textenglish[variant=american]{$E=\frac{\phi}{A}\left[\mathrm{\frac{W}{m^{2}}}\right]$}\tabularnewline
\hline 
שטף קרינה (של מקור קורן) & \textenglish[variant=american]{$\phi=IA_{sourc}\left[\mathrm{W}\right]$}\tabularnewline
\hline 
\begin{cellvarwidth}[t]
\raggedleft
עבור אנרגייה תמיד מתקיים

(כאשר ההספק קבוע)
\end{cellvarwidth} & \textenglish[variant=american]{$E=P\Delta t$}\tabularnewline
\hline 
ביטוי כללי להספק & \textenglish[variant=american]{$P=\frac{dE}{dt}\left[\mathrm{W}\right]$}\tabularnewline
\end{longtable}%
\end{minipage}}%
\fbox{\begin{minipage}[t]{0.5\columnwidth}%

\section*{האפקט הפוטו-אלקטרי}

\begin{longtable}[c]{r|r}
אנרגית פוטון כתלות בתדר & \textenglish[variant=american]{$E_{ph}=h\nu$}\tabularnewline
\hline 
\begin{cellvarwidth}[t]
\raggedleft
האנרגיה הקינטית של פוטון כאשר

$B$ אנרגיית הקשר של המתכת
\end{cellvarwidth} & \textenglish[variant=american]{$E_{k}=E_{ph}-B$}\tabularnewline
\hline 
תדירות הסף & \textenglish[variant=american]{$\nu=\frac{B}{h}\left[\mathrm{Hz}\right]$}\tabularnewline
\hline 
\begin{cellvarwidth}[t]
\raggedleft
\textbf{מתח עצירה} (כאשר $E_{k}=U_{e}$

האנרגייה הקינטית שווה לאנרגיה

הפוטנציאלית החשמלית)
\end{cellvarwidth} & \textenglish[variant=american]{$V_{stop}=\frac{E_{k}}{e}$}\tabularnewline
\hline 
\textenglish[variant=american]{} & \textenglish[variant=american]{$E=\frac{1240}{\lambda[\mathrm{nm}]}\left[\mathrm{eV}\right]$}\tabularnewline
\hline 
\textenglish[variant=american]{} & \textenglish[variant=american]{}\tabularnewline
\hline 
\textenglish[variant=american]{} & \textenglish[variant=american]{}\tabularnewline
\hline 
\textenglish[variant=american]{} & \textenglish[variant=american]{}\tabularnewline
\hline 
\textenglish[variant=american]{} & \textenglish[variant=american]{}\tabularnewline
\end{longtable}%
\end{minipage}}

\fbox{\begin{minipage}[t]{0.5\columnwidth}%
\textbf{שיווי משקל תרמודינמי}

\begin{longtable}[c]{>{\raggedleft}p{0.3\columnwidth}>{\raggedleft}p{0.6\columnwidth}}
\multicolumn{2}{r}{התפלגות בולצמן:~~~ $P(E)=Ae^{\frac{-E}{k_{B}T}};\,\,\ensuremath{k_{B}=1.38\times10^{-23}\left[\frac{J}{K}\right]}$}\tabularnewline
\textenglish[variant=american]{} & \textenglish[variant=american]{}\tabularnewline
תנאי נרמול: ~~~$\int P(E)dE=1$ & תוחלת של משתנה $x$:~~~$\left\langle x\right\rangle =\int^{\infty}_{-\infty}xP(x)dx$\tabularnewline
\multicolumn{2}{r}{ההסתברות למציאת חלקיק בגובה $h$:~~~$P(h)=\frac{mg}{k_{B}T}e^{\frac{-mgh}{k_{B}T}}$}\tabularnewline
\multicolumn{2}{r}{גובה ממוצע של חלקיק בעל מסה $m$:~~~$\left\langle h\right\rangle =\int hP(h)dh=\frac{k_{B}T}{mg}$}\tabularnewline
\multicolumn{2}{r}{האנרגיה הפוטנציאלית הממוצעת של חלקיק:~~~$\left\langle E_{p}\right\rangle =\frac{1}{2}m\omega^{2}\text{\ensuremath{\left\langle x^{2}\right\rangle }}=\frac{1}{2}k_{B}T$}\tabularnewline
\multicolumn{2}{r}{האנרגיה הקינטית הממוצעת של חלקיק:~~~$\left\langle E_{k}\right\rangle =\frac{1}{2}m\left\langle v^{2}\right\rangle =\frac{1}{2}k_{B}T$}\tabularnewline
\multicolumn{2}{>{\raggedleft}p{1\linewidth}}{עקרון החלוקה השווה: כל איבר בנוסחת האנרגיה הכוללת שתלוי במשתנה כלשהו
בריבוע, תורם בממוצע אנרגיה של \textenglish[variant=american]{$\frac{n}{2}k_{B}T$}
כאשר $n$ מספר דרגות החופש של החלקיק. דוג' לחלקיק חופשי במרחב 3 דרגות
חופש, תזוזה ב- $x,y,z$}\tabularnewline
\textenglish[variant=american]{} & \textenglish[variant=american]{}\tabularnewline
\textenglish[variant=american]{} & \textenglish[variant=american]{}\tabularnewline
\end{longtable}%
\end{minipage}}

{\fboxsep 5pt\fbox{\begin{minipage}[t]{0.5\columnwidth}%
\textbf{שיווי משקל תרמודינמי}

התפלגות בולצמן: $P(E)=Ae^{\frac{-E}{k_{B}T}}$%
\end{minipage}}}

\fbox{\begin{minipage}[t]{0.5\columnwidth}%
\begin{tabular}{cc}
\textbf{משוואת שרדינגר} & \textenglish[variant=american]{}\tabularnewline
המשוואה התלויה בזמן: & $\frac{\partial\Psi}{\partial t}=-\frac{\hbar^{2}}{2m}\frac{\partial^{2}\Psi}{\partial x^{2}}+U\Psi$$i\hbar$\tabularnewline
המשוואה שאינה תלויה בזמן: & $E\Psi(x)=-\frac{\hbar^{2}}{2m}\frac{\partial^{2}\Psi(x)}{\partial x^{2}}+U(x)\Psi(x)$\tabularnewline
\textenglish[variant=american]{} & \textenglish[variant=american]{}\tabularnewline
\end{tabular}%
\end{minipage}}

\fbox{\begin{minipage}[t]{0.5\columnwidth}%
$N_{n,\ell}$ קבוע הנרמול

המספר הקוונטי $\ell\ge\left|m_{\ell}\right|$ התנע הזוויתי של האלקטרון.
$m_{\ell}$ המספר הקוונטי-מגנטי של האלקטרון (ההיטל של $\ell$ על ציר
$\mathcal{Z}$) וקובע את כיוון הספין שלו.

תנאי קוונטיזציה%
\end{minipage}}

\end{hebrew}%
\begin{table}[H]
\centering
\begin{tabular}{|c|c|c|}
\hline 
\texthebrew{\textbf{משוואות שרדינגר לאטום המימן}} &  & \tabularnewline
\hline 
$\Psi(r,\theta,\varphi)=R\left(r\right)\Theta\left(\theta\right)\Phi\left(\varphi\right)$\texthebrew{\textbf{צורה
כללית}} &  & \tabularnewline
\hline 
 &  & \tabularnewline
\hline 
\begin{cellvarwidth}[t]
\centering
\texthebrew{%
\begin{minipage}[t]{0.15\columnwidth}%
\[
\kappa=\sqrt{\frac{2m_{e}|E|}{\hbar^{2}}}
\]
%
\end{minipage}}

\texthebrew{%
\begin{minipage}[t]{0.15\columnwidth}%
\[
\frac{1}{a_{0}}=\frac{m_{e}ke^{2}}{\hbar^{2}}
\]
%
\end{minipage}}

\texthebrew{~}
\end{cellvarwidth} & \begin{cellvarwidth}[t]
\centering
~

~

~

$\frac{\partial^{2}u(r)}{\partial r^{2}}+\left(-\kappa^{2}+\frac{2}{a_{0}r}-\frac{\ell(\ell+1)}{r^{2}}\right)u(r)=0$
\end{cellvarwidth} & \begin{cellvarwidth}[t]
\centering
~

~

~

\texthebrew{המשוואה הרדיאלית \textenglish[variant=american]{$R(r)$}}
\end{cellvarwidth}\tabularnewline
\hline 
 & $\frac{1}{\Theta(\theta)\sin\theta}\frac{\partial}{\partial\theta}\left(\sin\theta\frac{\partial\Theta(\theta)}{\partial\theta}\right)+\left(\ell(\ell+1)-\frac{m^{2}_{\ell}}{\sin^{2}\theta}\right)=0$ & \texthebrew{המשוואה הפולרית\textenglish[variant=american]{ $\Theta(\theta)$}}\tabularnewline
\hline 
$\Phi\left(\varphi\right)=\frac{1}{\sqrt{2\pi}}e^{im_{\ell}\varphi}$\texthebrew{פתרון:} & $\frac{\partial^{2}\Phi(\varphi)}{\partial\varphi^{2}}+m^{2}_{\ell}\Phi(\varphi)=0$ & \texthebrew{המשוואה האזימוטלית $\Phi(\varphi)$}\tabularnewline
\hline 
 &  & \tabularnewline
\hline 
\end{tabular}
\end{table}

\begin{hebrew}%
\textbf{אטום המימן}

\[
a_{0}=\frac{\hbar}{m_{e}ke^{2}}=0.53\left[\mathring{\mathrm{A}}\right]
\]

צורה כללית

פתרון המשוואה האזימוטלית:כאשר: 
\[
R_{n,\ell}\left(r\right)=N_{n,\ell}\left(\frac{2r}{a_{0}n}\right)^{\ell}e^{-\frac{r}{a_{0}n}}L^{2\ell+1}_{n+\ell}\left(\frac{2r}{a_{0}n}\right)
\]
$N_{n,\ell}$ קבוע הנרמול

המספר הקוונטי $\ell\ge\left|m_{\ell}\right|$ התנע הזוויתי של האלקטרון.
$m_{\ell}$ המספר הקוונטי-מגנטי של האלקטרון (ההיטל של $\ell$ על ציר
$\mathcal{Z}$) וקובע את כיוון הספין שלו.

תנאי קוונטיזציה

\[
0\le\ell\le n-1
\]
\end{hebrew}%

\end{document}
